\section{The problem, conventions, and scope}\label{bmh:sec:scope}

\subsection{A precise continuation of the repository}
The canonical manuscript \emph{Polylogarithms and their Arithmetic Bridges}
collects an extensive theory of polylogarithmic identities, harmonic sums,
Gamma arithmetic, Hurwitz jets, and Stieltjes correlations~\cite{Repo}.
The related report \emph{Mellin--Lerch Transforms and Dilation Identities for
Stieltjes Correlations} proves a one-factor Mellin transform and then asks,
in its tenth proposed research question, for a treatment of
\begin{equation}\label{bmh:eq:problem}
 \int_0^\infty\frac{x^{a-1}\Li_s(-zx)\Li_t(-wx)}{(1+x)^N}\dd x.
\end{equation}
Its two explicit objectives are finite multiple-polylogarithm closure at
integer Mellin parameters and a determination of the depth forced by
independent arguments~\cite{MLD}. The present article answers the first
objective when the polylogarithm orders are positive integers. It goes
further at equal arguments by deriving a finite double-Hurwitz formula
for noninteger $a$, exact formulas at both integer resonances in the
$N=1$ strip, and all logarithmic moments.

This is not a claim to settle the question for arbitrary complex spectral
orders, nor to prove numerical or motivic minimal depth. Generic
hyperlogarithmic integration is established mathematics; the multiscale
closure proof below is an explicit specialization of that method, with
its alphabet, full product convergence strip, and residue-weight
refinement made concrete~\cite{Panzer}. The connected double-Hurwitz
formula and the coefficient identities are derived here in detail and
are the main research contribution relative to the inspected repository
material. An exhaustive literature priority search has not been made.

The main source path is
\begin{center}\small
\url{Analysis/Polylogarithms/docs/reports/stieltjes-correlation-calculus/continuations/mellin-lerch-dilation/}.
\end{center}
The research-question file is \src{sections/06-audit-research.tex}; its
blob at the recorded pin is
\src{c0e3d58e42005ca844ca8b6900d7feaf1b33e2f8}.
The existing one-factor theorem, used with attribution in
Sections~\ref{bmh:sec:reciprocity}--\ref{bmh:sec:stieltjes}, is in
\src{sections/02-mellin.tex}. We also inspected the canonical README,
main source, selected depth and discovery sections, the incoming
inventory, and its intake instructions. The binary contents of the eight
archives in that inventory were not analytically audited. Consequently
our comparison is with the specified inspected text, not a claim of
exhaustive novelty against every incoming archive.

\subsection{Notation and normalization}
All unqualified $m,n,N$ are positive integers; $h,r$ are nonnegative
integers. Put $W=m+n$ and define
\begin{equation}\label{bmh:eq:M-definition}
 M_N(a;m,n;z,w)=\int_0^\infty
 \frac{x^{a-1}\Li_m(-zx)\Li_n(-wx)}{(1+x)^N}\dd x.
\end{equation}
We suppress $z,w$ when both are one. Powers of positive real numbers use
the real logarithm. The polylogarithm is on its principal branch; positive
$z,w$ therefore never put its argument on the positive-real branch cut.
More general scales are considered on
$\C\setminus(-\infty,0]$. Standard polylogarithm, Hurwitz-zeta, and Gamma
conventions are those of the DLMF~\cite{DLMFpoly,DLMFhurwitz,DLMFbeta}.

We use a semicolon in $\zeta(s;q)$ for the single Hurwitz parameter,
to distinguish it from the comma-separated multiple-zeta indices.
Multiple zeta values use \emph{decreasing} indices:
\begin{equation}\label{bmh:eq:MZV}
 \zeta(k_1,\ldots,k_d)=
 \sum_{n_1>\cdots>n_d\ge1}\frac1{n_1^{k_1}\cdots n_d^{k_d}},
 \qquad k_1\ge2,
\end{equation}
and the common-shift double Hurwitz value is
\begin{equation}\label{bmh:eq:Zq}
 Z_q(A,B)=\sum_{\nu>\mu\ge0}
 \frac1{(\nu+q)^A(\mu+q)^B},
 \qquad A\ge2,
 \quad B\ge1,\quad \Rea q>0.
\end{equation}
Thus $Z_1(A,B)=\zeta(A,B)$. When the spectral indices vary in a
proof, the same series is used wherever it is absolutely convergent. Each of these series, and any fixed number of
its $q$-derivatives, converges normally on compact subsets of
$\Rea q>0$. The inner sum is at most logarithmic when $B=1$ and bounded
when $B>1$; this proves the assertion directly. Spectral differentiation
in $B$ near one adds only logarithms and is also normally convergent.

We write $(s)_j=s(s+1)\cdots(s+j-1)$, with $(s)_0=1$, and
$\binom uv=0$ for integer $u\ge0$ and $v\notin\{0,\ldots,u\}$.
Generalized harmonic numbers are
$H_d^{(r)}=\sum_{j=1}^d j^{-r}$, with $H_0^{(r)}=0$.
The Stieltjes convention is
\begin{equation}\label{bmh:eq:Stieltjes-convention}
 \zeta(s;q)=\frac1{s-1}
 +\sum_{r=0}^\infty\frac{(-1)^r}{r!}\gamma_r(q)(s-1)^r,
 \qquad \gamma_0(q)=-\psi(q).
\end{equation}
An order derivative $\Li_s^{[r]}$ differentiates $s$, not the argument.

\subsection{What the principal results say}
Theorem~\ref{bmh:thm:master} is a finite formula at all positive integer
orders and noninteger Mellin parameters in the full product strip.
Theorem~\ref{bmh:thm:compact} gives an especially short unshifted double-zeta
answer. Theorems~\ref{bmh:thm:jets} and \ref{bmh:thm:kernel} give the complete
logarithmic and integer-kernel calculi. Theorem~\ref{bmh:thm:central}
identifies a parity sector with a single-zeta answer, and
Theorem~\ref{bmh:thm:Stieltjes} connects bilinear polylogarithmic integrals
to every generalized Stieltjes index. The algebraic checks are finite
implementation tests. The proofs of the infinite families are the
analytic and combinatorial arguments printed below.
